% \chapter{Introduction}
\chapter{INTRODUCTION}\label{ch:intro}
%\titlespacing\chapter{0pt}{0pt}{0pt}

Education relies heavily on assessment as a means of evaluating student learning, guiding curriculum development, and informing teaching practices. 
Among different assessment methods, written examinations stand out as a powerful tool for evaluating higher-order skills such as reasoning, problem-solving, and expressive ability. 
However, grading written exams poses significant challenges, particularly in contexts where linguistic complexity, large class sizes, and teacher workload combine to create barriers to fair and consistent evaluation. 
In Thailand, these challenges are even more pronounced due to the unique features of the Thai language and the increasing demand for efficiency in educational systems. This research addresses these issues 
by proposing an AI-assisted grading framework that leverages Large Language Models (LLMs) and rule-based reasoning to improve the evaluation of Thai written exam responses.

\section{Background and Motivation}
Examinations remain one of the most critical tools in educational assessment, 
serving as a benchmark for measuring students’ comprehension, knowledge retention, 
and ability to apply concepts. Among different types of examinations, written exams 
are particularly valuable because they provide a deeper insight into students’ higher-order 
cognitive abilities, such as critical reasoning, logical explanation, problem-solving, and 
language proficiency \citep{JungX2024}. Unlike multiple-choice questions that primarily measure recognition and recall, 
written responses reveal students’ thought processes, creativity, and capacity for structured argumentation.

Despite these advantages, written exam evaluation presents significant challenges. 
Traditional manual grading is time-consuming, resource-intensive, and often inconsistent. 
Human evaluators, including teachers and teaching assistants (TAs), are susceptible to fatigue, 
emotional states, and subjective interpretation, which may lead to inconsistency in grading outcomes \citep{Mansson2016}. 
Even when grading rubrics are provided, variations in understanding and application among 
TAs may result in uneven assessment standards.

% \cite{Garlotta2001} stated that “Sirindhorn International Institute of Technology
% (SIIT) offers three master of engineering programs, namely, Master of Engineering
% Program in Engineering Technology, Master of Engineering Program in Information and 
% Communication Technology for Embedded Systems, and Master of Engineering 
% Program in Logistics and Supply Chain Systems Engineering” (Section~\ref{ss:ref}).

\section{Challenges in the Thai Language Context}

In the Thai educational context, the grading of written responses is further 
complicated by the unique linguistic and structural properties of the Thai 
language. Thai grammar is highly flexible, allowing for multiple valid 
sentence constructions. The language is tonal, meaning word meaning can 
change depending on pitch, and it contains numerous context-dependent 
exceptions \citep{Tongsilp2024}. Moreover, Thai writing omits spaces between words, which 
complicates segmentation and tokenization processes required for 
computational analysis.

These linguistic features make it particularly difficult for evaluators to 
interpret student responses accurately, especially when assessing reasoning, 
argumentation, or nuanced explanations. Teachers must balance linguistic 
expertise with an understanding of the subject matter and grading criteria. 
Given the open-ended nature of written responses, the risk of 
misinterpretation and grading bias is substantially higher compared to 
multiple-choice formats.

As a result, many Thai educators prefer multiple-choice examinations, which 
are quicker and easier to grade. This choice, however, comes at the expense 
of assessing higher-order skills. On average, teachers in Thailand manage 
classes of 100–400 students per subject while teaching 7–21 hours weekly 
and handling administrative, extracurricular, and pastoral 
responsibilities \citep{Goksoy2014}. Such workloads leave little room for the detailed 
evaluation of written work, leading to an overreliance on more easily 
graded exam types.

% The study performed by \cite{Tokiwa2006} showed that Sirindhorn
% International Institute of Technology (SIIT) offers three master of engineering
% programs, namely, Master of Engineering Program in Engineering Technology, Master 
% of Engineering Program in Information and Communication Technology for Embedded 
% Systems, and Master of Engineering Program in Logistics and Supply Chain Systems
% Engineering \citep{Bocchini2013, Montanari2003}. 

\section{Limitations of Current Assessment Practices}

Although multiple-choice exams are efficient, they fail to capture the 
richness of student reasoning, creativity, and expressive ability 
\citep{scouller1998influence, birenbaum1997alternative}. 
Written exams, despite their resource-intensive nature, remain the most 
effective tool for evaluating comprehensive understanding and critical 
thinking \citep{gipps2011beyond, gierl2017automated}. 
This creates a tension between educational goals and practical 
constraints.

In recent years, efforts to explore the potential of LLMs in grading have included rubrics, 
peer evaluation, and partial automation through keyword-based scoring 
systems \citep{lu2021automatic, ridley2022automated}. 
However, these approaches often fall short due to the complexities 
of natural language and the diversity of student expression. Rule-based 
systems, for instance, can detect certain patterns but struggle with 
semantic nuances, paraphrasing, and context-dependent meaning 
\citep{dikli2006automated, shermis2014state}.

\section{Role of Bloom’s Taxonomy in Assessment}

Bloom’s Taxonomy provides a structured framework for categorizing cognitive skills, 
ranging from basic recall of facts to higher-order reasoning and critical thinking 
\citep{bloom1956taxonomy,anderson2001taxonomy}. 
In the context of this research, Bloom’s Taxonomy serves two important purposes. 

First, it provides a pedagogical lens for understanding the type of cognitive 
process each exam question is designed to assess. This ensures that grading 
criteria align with intended learning outcomes \citep{krathwohl2002revision}. 

Second, Bloom’s levels are used as a pre-screening mechanism for the 
AI-assisted framework. Each exam question is labelled with a Bloom’s 
category by expert teachers, and this label is then used to guide 
model selection. In practice, different types of questions 
(e.g., recall vs. analysis) may be better handled by different LLMs or 
combinations of LLM and rule-based reasoning \citep{zohar2004higher,stanny2016reexamining}. 

By integrating Bloom’s Taxonomy into the framework, the system 
acknowledges both pedagogical validity and technical feasibility, 
linking educational theory with AI-based proof-of-concept evaluation.


\section{Advances in Natural Language Processing (NLP)}

Recent breakthroughs in Natural Language Processing (NLP) and the advent 
of Large Language Models (LLMs) have opened new possibilities for automated 
exam evaluation \citep{Agarwal2021}. LLMs demonstrate strong capabilities in tasks such as 
text classification, summarization, sentiment analysis, and automated 
evaluation across multiple languages. Unlike traditional rule-based methods, 
LLMs leverage deep learning and attention mechanisms to capture contextual 
meaning, making them suitable for handling open-ended and nuanced responses.

For the Thai language, the application of LLMs introduces unique 
opportunities and challenges. While most existing LLMs are trained 
primarily on English or other widely spoken languages, recent developments 
have included multilingual models capable of handling Thai to varying 
degrees. By integrating LLMs with rule-based reasoning, it may be possible 
to design a hybrid evaluation framework that combines the flexibility and 
contextual understanding of LLMs with the consistency and transparency of 
structured rules.

\section{Research Motivation and Contribution}

The motivation behind this research is twofold:
\begin{enumerate}
    \item To reduce the grading burden on Thai teachers while improving fairness and consistency in written exam evaluation.
    \item To promote the wider adoption of open-ended questions in Thai education, thereby encouraging the development of critical reasoning, creativity, and expressive skills in students.
\end{enumerate}

This thesis aims to develop an AI-assisted grading framework for 
Thai-language written exam questions. By leveraging LLMs in combination 
with rule-based reasoning, the framework is designed to enhance accuracy, 
consistency, and efficiency in grading, while maintaining interpretability 
for educators.

The expected contributions of this research are as follows:
\begin{itemize}
    \item A comprehensive analysis of challenges in grading Thai written responses.
    \item A proof-of-concept framework combining LLMs and rule-based reasoning.
    \item Empirical evaluation of the framework’s performance compared to human graders.
    \item Insights into the broader implications for educational assessment and curriculum design.
\end{itemize}

\section{Research Objectives}

The objectives of this research are as follows:
\begin{enumerate}
    \item To investigate the challenges and limitations in grading Thai written exam responses.
    \item To design a proof-of-concept AI-assisted grading framework that integrates LLMs with rule-based reasoning.
    \item To evaluate the framework’s effectiveness against human grading in terms of fairness, consistency, and interpretability.
    \item To provide insights into how such frameworks may support, rather than replace, teachers in Thai educational contexts.
\end{enumerate}

\section{Scope of the Study}

This study is limited to the following scope:
\begin{enumerate}
    \item Focuses on Thai-language written exam questions at the secondary school level.
    \item Examines only short-answer and essay-type responses (0–3 scoring scale).
    \item Uses a proof-of-concept approach with selected LLMs and teacher-designed rubrics.
    \item Each exam question is pre-classified into Bloom’s Taxonomy levels by expert teachers, and these labels are used to guide model selection within the framework.
    \item The framework is evaluated on fairness, consistency, and interpretability, but not on large-scale deployment or integration into actual school systems.
\end{enumerate}

\section{Structure of the Thesis}

The remainder of this thesis is organized as follows:
\begin{itemize}
    \item \textbf{Chapter 2 – Literature Review}: Reviews existing work on exam evaluation, NLP in education, and applications of LLMs in AI-assisted grading frameworks.
    \item \textbf{Chapter 3 – Methodology}: Describes the proposed framework, including data collection, preprocessing, model selection, and rule-based integration.
    \item \textbf{Chapter 4 – Result and Discussion}: Presents the experimental findings, discusses the strengths and limitations of the framework, and compares it with baseline methods.
    \item \textbf{Chapter 5 – Conclusion and Future Work}: Summarizes key findings, highlights contributions, and suggests directions for future research.
\end{itemize}